\providecommand{\SysName}{CARE}

\section{Implementation}
\label{sec:implementation}

\noindent\textbf{Prototype size and interface.}
We implemented \SysName{} as a Python research prototype for C/C++ codebases.
At the time of writing, the implementation contains 14,264 lines of Python code
in the core package and experiment scripts, excluding 2,576 lines of pytest
tests. The command-line entry point is \texttt{care run}, which accepts a target
project, optional build and test commands, an optional compile database,
analysis backends, detector-confirmation options, LLM mode, and repair-loop
parameters. Runtime options are normalized through a \texttt{CareConfig}
dataclass, and all cross-stage artifacts, including functions, graph edges,
resource events, opportunities, plans, patch candidates, and validation results,
are JSON-serializable dataclasses.

\noindent\textbf{Analysis and detection.}
\SysName{} supports both precise and scalable analysis paths. When clang and
compile arguments are available, it extracts function facts from clang JSON AST
dumps and can use clang analyzer CFG output. Otherwise, it falls back to a
deterministic brace-aware parser that masks comments and strings, identifies
function bodies, and extracts calls, returns, labels, gotos, and local facts.
The context builder merges these facts with lightweight CFGs, call edges,
data-flow summaries, resource events, side-effect summaries, and security
patterns. The detector orchestrator runs goto, resource, duplicate-check,
dead-code, and security-smell detectors, then deduplicates, clusters, filters,
optionally confirms high-impact findings with clang, and ranks opportunities by
severity and evidence confidence.

\noindent\textbf{LLM integration and patch materialization.}
The LLM layer is deliberately small: a generic client exposes one completion
method, with a deterministic mock provider for tests and an OpenAI-compatible
provider configured only through environment variables. The planner converts an
opportunity into a structured safety contract containing affected files,
affected functions, invariants, security constraints, behavior-preservation
goals, and validation strategy. The patch generator asks for either a JSON edit
plan or a unified diff. JSON edit plans are converted into diffs using the
current project files; raw diffs are normalized and checked for valid unified
diff structure. Before validation, \SysName{} repairs common LLM path mistakes,
such as duplicated project prefixes, missing \texttt{a/}/\texttt{b/} prefixes,
and absolute paths that still point inside the target project.

\noindent\textbf{Validation, reporting, and experiments.}
Every patch candidate is validated from a clean project state using git when
available and a temporary backup otherwise. The validator checks patch
application, rejects no-effect patches, runs the configured build and test
commands, invokes optional static analyzers when installed, and applies
resource, semantic, and security regression checkers. Failures are classified
into a taxonomy covering patch-application errors, compile/test failures,
resource regressions, semantic regressions, security regressions, and LLM/API
failures; the repair loop feeds these logs back into smaller repair prompts.
Reports are emitted as JSON, Markdown, and per-candidate diff files. For the
large-scale OSS50 experiments, we implemented resumable queue scripts, a
project-sharded process runner, and a dynamic work-stealing runner that validates
remaining opportunities in isolated per-worker project copies.
